% !TeX root = ../../main.tex
\chapter{总结与展望}\label{chap:conclusion}
\section{全文研究工作总结}
本文面向以 SpMV 为核心的稀疏矩阵计算流程，研究如何利用不同计算阶段、数据区域及运行状态的需求差异，通过近似计算及软硬件协同设计降低实际执行成本。针对固定精度与固定范围执行可能引入的冗余，全文依次研究低比特分层计算、结构感知选择性计算和状态感知混合精度计算。

首先，本文从 Top-K 检索与 CG 迭代中提取共同的系统问题，分析稀疏数据的低算术强度、不规则访问及并行负载差异。通过已有实验分别说明粗筛与重算的精度差异、候选区域的贡献差异及迭代状态的数值敏感性。三者均依据计算需求调整实际工作量，并分别以候选覆盖、检索召回和数值收敛约束近似过程，形成共同的系统设计思路与各自独立的质量评价体系。

其次，针对全量高精度计算成本，AccelES 采用 INT6 预筛与 FP32 候选重算，结合 ReSparse 降低低比特阶段的数据规模。Ultra-CSR 与 Random-CSR 分别适配流式扫描和候选按行读取，多 Q-core 及共享 F-core 匹配两阶段工作量。实验中完整重算配置相对所选 FPGA、GPU 及 CPU 基线分别取得3.4、2.5 及153.3 倍加速，同时最大候选保持率损失为0.95\%。该工作表明，高精度运算可以依据输出需求集中到少量候选，但粗筛遗漏仍构成质量边界。

再次，针对低比特阶段仍需全量扫描的问题，ParetoES 采用方向敏感聚类和簇级探测，缩减实际计算范围。ACPE 将簇筛选、地址定位和簇内计算统一在流式执行中，分布式微型排序降低比较逻辑规模，\hbalance{}缓解查询热点导致的核间执行拖尾。在原论文 Recall@100 为0.8 至1.0 的评估区间内，相对 AccelES 平均吞吐提升为2.27 倍。接近全召回时收益减小，说明选择性计算的优势依赖数据结构及质量目标。

最后，针对迭代状态相关的精度需求，EvoCG 采用强化学习控制局部精度，并通过 Tile-CSC、统一 SpMV 内核与低开销运行时执行异构动作。默认矩阵专属策略在47 个测试矩阵上完成求解，相对 cuSPARSE 等系统取得端到端改进；与相同执行基础的 FP64 配置相比，精度控制贡献1.28 倍几何平均加速。策略控制平均占在线时间1.09\%，但训练和多精度存储需要额外成本，跨矩阵迁移也并非全部成功。

表\ref{tab:summary}将三项工作放在同一分析框架下总结。比较围绕研究对象、决策机制和适用条件展开，各项性能结果分别对应其硬件平台、数据集与质量目标。
\begin{table}[htbp]\centering
\caption{三项研究的决策机制、协同实现与适用条件}\label{tab:summary}
\begin{tabularx}{\textwidth}{p{2cm}XXX}\toprule
比较项目&AccelES&ParetoES&EvoCG\\\midrule
主要问题&高精度全量计算成本&低比特全量候选访问&状态相关的精度选择\\
决策依据&近似得分与元素幅值&查询与候选簇结构&当前搜索方向数值分布\\
近似机制&低比特筛选、元素剪枝&簇级候选范围缩减&局部动态混合精度\\
执行平台&HBM FPGA&HBM FPGA&A100 GPU\\
关键协同&双格式、异构两阶段&ACPE、微排序、热点均衡&Tile-CSC、统一内核、运行时\\
主要代价&双副本、候选遗漏与重算&索引构建、簇跳转与召回损失&训练、容量增长与额外迭代\\
关系定位&降低单位计算成本&进一步缩减计算范围&拓展到反馈驱动的数值迭代\\\bottomrule
\end{tabularx}
\end{table}

\section{本文主要创新点}
（1）提出面向阶段精度差异的低比特分层计算与异构执行方法。将 Top-K 候选识别和候选得分生成分离，以低比特运算和非零剪枝减少全量高精度访问，设计面向两种访问模式的数据格式及 FPGA 计算路径。通过近似筛选、高精度重算和数据流结构的协同，使精度预算与阶段计算需求相匹配。

（2）提出面向候选贡献差异的结构感知选择性计算与多核执行方法。在前一工作低比特流式基础上，采用稀疏方向敏感聚类与可配置簇探测，将召回需求映射到实际访问范围；以 ACPE、分布式微型排序和热点均衡支撑不规则选择性执行。在继承 INT6 与 Ultra-CSR 的基础上，实现候选空间优化与多核执行机制的协同。

（3）提出面向运行状态的闭环混合精度控制与 GPU 协同执行方法。以搜索方向片段表征数值状态，通过强化学习决定局部精度，并结合 Tile-CSC 共享结构索引、线程束级规则执行及低开销迭代运行时，将精度变化转化为求解时间改善。该工作明确区分精度控制与执行基础的收益，分析数值质量、策略迁移及训练摊销边界。

三项研究共同支持以下系统设计认识：近似计算提供减少数据访问与运算的机会，但端到端收益还取决于质量约束、决策代价和执行效率。只有使近似粒度与数据表示、并行资源及运行时机制相匹配，才可能稳定实现这一收益。本文的实验与分析从三类计算过程验证了这一设计原则，并揭示了其适用条件。

\section{研究局限与未来工作展望}
\subsection{质量约束与可靠近似机制}
本文检索方案主要通过验证查询校准候选保持与召回，尚未建立对任意查询均适用的遗漏界。后续可结合候选分数间隔、量化误差界和在线置信度，动态扩大候选或探测范围，使固定预算转变为质量驱动预算。需要同时控制质量检查开销，避免把减少的扫描重新消耗在复杂判定中。

CG 方面，学习策略在有限矩阵集合上表现出较好效果，但缺少适用于任意矩阵、右端项和预条件器的收敛保证。后续可研究将可验证数值准则与学习策略结合，在必要时触发残差重算、精度提升或可靠重启。此类机制应重新分析额外同步与验证成本，不能直接假定安全回退没有性能代价。

\subsection{动态数据、查询分布与泛化能力}
AccelES 和 ParetoES 均依赖可复用的离线准备，当前实验不能完整代表频繁插入、删除或向量更新的数据库。可进一步研究增量量化、局部重聚类和轻量数据迁移，并以持续更新条件下的吞吐与尾延迟评价，而非只测试静态索引。

\hbalance{}使用代表查询估计热点，分布偏移会降低映射效果。后续可结合在线统计与更新收益预测，决定何时重新平衡。EvoCG 的跨矩阵迁移则显示数值分布具有一定共享性但并不普适，可进一步研究任务族划分、多矩阵训练与策略有效性检测，以减少专属训练次数，同时明确迁移失败的处理成本。

\subsection{数据表示与硬件执行的进一步协同}
本文 FPGA 设计围绕较低维度稀疏嵌入与固定 HBM 接口展开。更高维度和更大 $K$ 需要调整索引字段、查询缓存和候选维护结构。未来可研究可变长索引、分层查询缓存及按数据分布调节候选资源的方法，并考察这种灵活性带来的控制开销。

Tile-CSC 通过多份数值流换取低成本精度切换，在部分极分散矩阵上存在较高存储开销。可进一步研究按需生成数值流、分层压缩或混合布局，降低容量压力。相关设计必须保留统一遍历路径的优势，避免节省存储后引入大量在线转换和分支。

\subsection{更完整的端到端评估}
当前三项工作使用不同平台和工作负载，尚不能形成统一环境下的综合性能排序。后续可建立涵盖质量目标、预处理摊销、更新频率和能耗范围的共同评估规范。对检索系统，应区分候选检索质量与下游任务效果；对数值求解，应同时报告停止判据、真实残差、失败次数及训练成本。

此外，单板单设备结果尚未覆盖多设备通信和更大规模数据分布。跨设备扩展时，近似策略可能减少本地计算却增加不均衡通信，局部控制也可能与全局同步耦合。未来应在保持计算语义的前提下，研究近似决策与通信划分、全局调度之间的关系，并用完整应用时间而非单一算子吞吐评价其价值。
